%!TEX program = xelatex
%!TEX root = ../all_beamer.tex

%% Decision Trees – Beamer Presentation
%% Based on notes by Giorgio Gambosi, University of Rome "Tor Vergata"
%% Style modelled on ml-mlp-beamer.tex

%% !TEX root = all_beamer.tex
% ==============================================================================
% HEADER LATEX PER PRESENTAZIONI BEAMER
% Corso di Machine Learning - Giorgio Gambosi
% Università di Roma "Tor Vergata"
% ==============================================================================

% ------------------------------------------------------------------------------
% CLASSE DOCUMENTO E METADATI
% ------------------------------------------------------------------------------
\documentclass[aspectratio=169, 10pt]{beamer}

\subtitle{Course of Machine Learning}
\author{Giorgio Gambosi}
\institute[Tor Vergata]{Master's Degree in Computer Science\\University of Rome ``Tor Vergata''}
\date{A.Y. 2025--2026}

% ------------------------------------------------------------------------------
% PACCHETTI MATEMATICI
% ------------------------------------------------------------------------------
\usepackage{amsmath, amssymb, amsfonts, mathtools, amsthm}
\usepackage{bm}              % Simboli matematici in grassetto
\usepackage{upgreek}         % Lettere greche dritte
\usepackage{derivative}      % Gestione semplificata delle derivate

% ------------------------------------------------------------------------------
% PACCHETTI GRAFICI E VISUALIZZAZIONE
% ------------------------------------------------------------------------------
\usepackage{graphicx}        % Inclusione immagini
\usepackage{xcolor}          % Gestione colori
\usepackage{xkcdcolors}      % Colori aggiuntivi stile xkcd
\usepackage{microtype}       % Miglioramenti tipografici

% ------------------------------------------------------------------------------
% PACCHETTI PER TABELLE
% ------------------------------------------------------------------------------
\usepackage{booktabs}        % Tabelle professionali
\usepackage{array}           % Estensioni per tabelle
\usepackage{multirow}        % Celle multi-riga
\usepackage{multicol}        % Layout multi-colonna

% ------------------------------------------------------------------------------
% TIKZ E PGFPLOTS
% ------------------------------------------------------------------------------
\usepackage{tikz}
\usepackage{tikz-network}
\usepackage[customcolors]{hf-tikz}
\usepackage{pgfplots}
\pgfplotsset{compat=1.18}
\usepgfplotslibrary{fillbetween}

% Librerie TikZ
\usetikzlibrary{arrows, arrows.meta, decorations.pathmorphing, decorations.pathreplacing}
\usetikzlibrary{decorations.markings, snakes, positioning, backgrounds, fit, shadows}
\usetikzlibrary{automata, patterns, calc, intersections, through, shapes}
\usetikzlibrary{shapes.geometric, shapes.arrows, scopes, chains, matrix}

% ------------------------------------------------------------------------------
% ALTRI PACCHETTI
% ------------------------------------------------------------------------------
\usepackage[skins]{tcolorbox}  % Box colorati
\usepackage{subcaption}         % Sottocaption per figure
\usepackage{listings}           % Inclusione codice sorgente

% ------------------------------------------------------------------------------
% TEMA E STILE BEAMER
% ------------------------------------------------------------------------------
\usetheme{Madrid}
\usecolortheme{whale}
\useinnertheme{rounded}
\usefonttheme{professionalfonts}
\setbeamertemplate{navigation symbols}{}  % Rimuove simboli di navigazione

% ------------------------------------------------------------------------------
% DEFINIZIONE COLORI PERSONALIZZATI
% ------------------------------------------------------------------------------
% Colori principali del tema
\definecolor{MainBlue}{RGB}{31,56,100}      % Blu principale (header, titoli)
\definecolor{AccentBlue}{RGB}{70,130,180}   % Blu accento (blocchi)
\definecolor{AlertRed}{RGB}{160,30,30}      % Rosso per alert
\definecolor{ExGreen}{RGB}{30,120,60}       % Verde per esempi
\definecolor{BlockBg}{RGB}{220,232,246}     % Sfondo blocchi
\definecolor{torgray}{RGB}{220,225,235}     % Grigio Tor Vergata

% Colori alternativi (Tor Vergata style)
\definecolor{TvDark}{RGB}{30,50,80}         % Blu scuro header
\definecolor{TvMid}{RGB}{70,100,140}        % Blu medio blocchi
\definecolor{TvLight}{RGB}{210,220,235}     % Blu chiaro sfondi
\definecolor{TvAlert}{RGB}{160,30,30}       % Rosso alert
\definecolor{TvGreen}{RGB}{20,100,50}       % Verde esempi

% Colori supplementari
\definecolor{torred}{RGB}{153,0,0}
\definecolor{torcyan}{RGB}{0,112,192}
\definecolor{cerise}{HTML}{DA3A62}
\definecolor{teal}{HTML}{029386}
\definecolor{slateblue}{HTML}{5B5EA6}
\definecolor{orange}{HTML}{F97306}

% Colori per nodi TikZ
\definecolor{nodeblue}{RGB}{101,140,196}    % xkcd Faded Blue

% ------------------------------------------------------------------------------
% CONFIGURAZIONE COLORI BEAMER
% ------------------------------------------------------------------------------
\setbeamercolor{frametitle}{bg=MainBlue, fg=white}
\setbeamercolor{title}{bg=MainBlue, fg=white}
\setbeamercolor{structure}{fg=MainBlue}

% Blocchi standard
\setbeamercolor{block title}{bg=AccentBlue, fg=white}
\setbeamercolor{block body}{bg=BlockBg}

% Blocchi alert
\setbeamercolor{block title alerted}{bg=AlertRed, fg=white}
\setbeamercolor{block body alerted}{bg=red!8}

% Blocchi esempio
\setbeamercolor{block title example}{bg=ExGreen, fg=white}
\setbeamercolor{block body example}{bg=green!7}

% ------------------------------------------------------------------------------
% CONFIGURAZIONE FONT BEAMER
% ------------------------------------------------------------------------------
\setbeamerfont{frametitle}{size=\normalsize,series=\bfseries}
\setbeamerfont{block title}{size=\small,series=\bfseries}

% ------------------------------------------------------------------------------
% STILI TIKZ PERSONALIZZATI
% ------------------------------------------------------------------------------
\tikzset{
  >=latex,  % Frecce stile LaTeX di default
  % Nodi per reti neurali
  node/.style={thick,circle,draw=myblue,minimum size=22,inner sep=0.5,outer sep=0.6},
  node in/.style={node,green!20!black,draw=mygreen!30!black,fill=mygreen!25},
  node hidden/.style={node,blue!20!black,draw=myblue!30!black,fill=myblue!20},
  node convol/.style={node,orange!20!black,draw=myorange!30!black,fill=myorange!20},
  node out/.style={node,red!20!black,draw=myred!30!black,fill=myred!20},
  % Connessioni
  connect/.style={thick,mydarkblue},
  connect arrow/.style={-{Latex[length=4,width=3.5]},thick,mydarkblue,shorten <=0.5,shorten >=1},
  % Stili numerati per facile mappatura
  node 1/.style={node in},
  node 2/.style={node hidden},
  node 3/.style={node out}
}

% ------------------------------------------------------------------------------
% AMBIENTI TEOREMI
% ------------------------------------------------------------------------------
\setbeamertemplate{theorems}[numbered]
\newtheorem{mytheorem}{Theorem}
\newtheorem{mydefinition}{Definition}

% ------------------------------------------------------------------------------
% CONFIGURAZIONE LISTINGS (CODICE)
% ------------------------------------------------------------------------------
\lstset{
  language=Python,
  basicstyle=\ttfamily\tiny,
  keywordstyle=\color{MainBlue}\bfseries,
  commentstyle=\color{gray},
  frame=single,
  breaklines=true,
}

% ==============================================================================
% MACRO MATEMATICHE PERSONALIZZATE
% ==============================================================================

% ------------------------------------------------------------------------------
% SPAZI CALLIGRAFICI (Calligraphic spaces)
% ------------------------------------------------------------------------------
\providecommand{\calX}{\mathcal{X}}         % Spazio input
\providecommand{\calY}{\mathcal{Y}}         % Spazio output
\providecommand{\calH}{\mathcal{H}}         % Spazio ipotesi
\providecommand{\calT}{\mathcal{T}}         % Training set
\providecommand{\calA}{\mathcal{A}}         % Algoritmo
\providecommand{\calR}{\mathcal{R}}         % Rischio
\providecommand{\calF}{\mathcal{F}}         % Famiglia di funzioni
\providecommand{\calL}{\mathcal{L}}         % Loss function

% Alias per compatibilità
\providecommand{\Rr}{\mathcal{R}}
\providecommand{\Tr}{\mathcal{T}}
\providecommand{\Hr}{\mathcal{H}}

% ------------------------------------------------------------------------------
% RISCHIO E PROBABILITÀ
% ------------------------------------------------------------------------------
\providecommand{\barR}{\overline{\mathcal{R}}}  % Rischio empirico
\providecommand{\EmpR}{\overline{\mathcal{R}}}  % Rischio empirico (alias)
\providecommand{\Rbar}{\overline{\mathcal{R}}}  % Rischio empirico (alias)
\providecommand{\pM}{p_M}                       % Probabilità modello
\providecommand{\pC}{p_C}                       % Probabilità corretta
\providecommand{\Prob}[2]{\mathbb{P}_{#1}\!\left[#2\right]}  % Probabilità

% ------------------------------------------------------------------------------
% FUNZIONI INDICATRICI
% ------------------------------------------------------------------------------
\providecommand{\indicator}[1]{\mathbf{1}\!\left[#1\right]}  % Funzione indicatrice
\providecommand{\ind}[1]{\mathbf{1}[#1]}                      % Versione breve

% ------------------------------------------------------------------------------
% OPERATORI MATEMATICI
% ------------------------------------------------------------------------------
\providecommand{\argmax}{\operatorname*{arg\,max}}  % Argomento del massimo
\providecommand{\argmin}{\operatorname*{arg\,min}}  % Argomento del minimo
\providecommand{\vcd}[1]{\operatorname{VCdim}(#1)} % VC dimension
\providecommand{\defeq}{\stackrel{\text{def}}{=}}  % Definito come

% ------------------------------------------------------------------------------
% VETTORI E MATRICI (Bold Math)
% ------------------------------------------------------------------------------
% Vettori minuscoli
\providecommand{\bx}{\mathbf{x}}            % Vettore x
\providecommand{\bz}{\mathbf{z}}            % Vettore z
\providecommand{\bw}{\mathbf{w}}            % Vettore peso
\providecommand{\btt}{\mathbf{t}}           % Vettore target
\providecommand{\wvec}{\mathbf{w}}          % Vettore peso (alias)
\providecommand{\bth}{\boldsymbol{\theta}}  % Parametro theta

% Matrici maiuscole
\providecommand{\bX}{\mathbf{X}}            % Matrice dati
\providecommand{\bZ}{\mathbf{Z}}            % Matrice Z
\providecommand{\Xmat}{\mathbf{X}}          % Matrice dati (alias)
\providecommand{\Hmat}{\mathbf{H}}          % Matrice H
\providecommand{\Smat}{\mathbf{S}}          % Matrice scatter generica
\providecommand{\bS}{\mathbf{S}}            % Matrice scatter
\providecommand{\Ii}{\mathbf{I}}            % Matrice identità
\providecommand{\bZero}{\mathbf{0}}         % Vettore zero
\providecommand{\bPsi}{\boldsymbol{\Psi}}   % Matrice Psi

% Medie (overline)
\providecommand{\ox}{\overline{\mathbf{x}}} % Media vettore x
\providecommand{\ow}{\overline{\mathbf{w}}} % Media vettore w
\providecommand{\oX}{\overline{\mathbf{X}}} % Media matrice X

% Matrici scatter specifiche (analisi discriminante)
\providecommand{\SW}{\mathbf{S}_W}          % Within-class scatter
\providecommand{\SB}{\mathbf{S}_B}          % Between-class scatter
\providecommand{\ST}{\mathbf{S}_T}          % Total scatter
\providecommand{\GX}{G_{\mathbf{X}}}        % Matrice Gram

% ------------------------------------------------------------------------------
% COMANDI DI FORMATTAZIONE TESTO
% ------------------------------------------------------------------------------
\providecommand{\term}[1]{\textbf{\color{accent}#1}}        % Termine importante
\providecommand{\halert}[1]{{\color{AlertRed}\textbf{#1}}}  % Alert rosso
\providecommand{\mlalert}[1]{{\color{AccentBlue}\textbf{#1}}}  % Alert blu
\providecommand{\mlemph}[1]{{\color{MainBlue}\textit{#1}}}  % Enfasi blu

% ==============================================================================
% FINE HEADER
% ==============================================================================

\newcommand{\blankpage}{
\clearpage
\thispagestyle{empty}
\null
\clearpage
}

%
%%% ── Packages ────────────────────────────────────────────────────────────────
%\usepackage{amsmath,amssymb,bm}
%\usepackage{booktabs}
%\usepackage{graphicx}
%\usepackage{xcolor}
%\usepackage{xkcdcolors}
%\usepackage{tikz}
%\usepackage{subcaption}

%% ── Math shorthands ─────────────────────────────────────────────────────────
%\newcommand{\SW}{\mathbf{S}_W}
%\newcommand{\SB}{\mathbf{S}_B}
%\newcommand{\ST}{\mathbf{S}_T}
%\newcommand{\Smat}{\mathbf{S}}
%\newcommand{\GX}{G_{\mathbf{X}}}

%% ── Metadata ────────────────────────────────────────────────────────────────
\title[Decision Trees]{Decision Trees}
%\subtitle{Course of Machine Learning}
%\author{Giorgio Gambosi}
%\institute[Tor Vergata]{Master's Degree in Computer Science\\University of Rome ``Tor Vergata''}
%\date{A.Y. 2025--2026}

%% ════════════════════════════════════════════════════════════════════════════
%\begin{document}

\begin{frame}
  \titlepage
\end{frame}

\begin{frame}{Outline}
  \tableofcontents
\end{frame}

%% ════════════════════════════════════════════════════════════════════════════
\section{Introduction to Decision Trees}

\begin{frame}{What Is a Decision Tree?}
  \framesubtitle{A supervised model based on recursive partitioning}
  \begin{itemize}
    \item A \alert{decision tree} is a supervised learning model for classification or regression.
    \item Core idea: build a sequence of \emph{if--then} decisions that progressively refine
          knowledge about the output.
    \item Starting from the full dataset, the tree \alert{repeatedly partitions} the instance space
          into smaller, more homogeneous regions.
    \item Formally: a \alert{recursive partition} of the feature space, where each internal decision
          reduces ambiguity and each terminal region provides a final prediction.
  \end{itemize}
  \begin{block}{Key appeal}
    Decision trees combine \alert{predictive power} with \alert{interpretability}: each internal
    node corresponds to a human-readable test on one input feature, and a prediction can be
    explained by tracing the path from root to leaf.
  \end{block}
\end{frame}

\begin{frame}{Why Decision Trees?}
  \begin{itemize}
    \item Unlike many machine learning models, a decision tree does \alert{not hide its logic}
          inside a large collection of numerical parameters.
    \item Its reasoning is laid out openly as a \alert{tree of conditions} that any practitioner
          can inspect, question, and communicate to a non-technical audience.
    \item Transparency is a decisive practical advantage in domains such as:
      \begin{itemize}
        \item Medicine (clinical decision support).
        \item Law (rule-based decisions).
        \item Finance (credit scoring, fraud detection).
      \end{itemize}
    \item In such settings, the \alert{rationale} behind a prediction often matters as much as
          the prediction itself.
  \end{itemize}
\end{frame}

%% ════════════════════════════════════════════════════════════════════════════
\section{Structure and Components}

\begin{frame}{Tree Structure: Nodes and Edges}
  \framesubtitle{Rooted hierarchical structure}
  \begin{columns}[t]
    \column{0.48\textwidth}
      \begin{block}{Root node}
        Sits at the top; represents the \alert{entire training set} before any partition.
        Starting point from which all decisions flow.
      \end{block}
      \vspace{0.4em}
      \begin{block}{Internal nodes}
        Each implements a \alert{decision rule} based on a feature. Routes an instance to
        one of its children depending on the test outcome.
      \end{block}
    \column{0.48\textwidth}
      \begin{block}{Leaf nodes (terminal nodes)}
        Represent regions where \alert{no further splitting} is performed. A final
        prediction is issued here: a class label, class distribution, or numerical value.
      \end{block}
      \vspace{0.4em}
      \begin{alertblock}{Traversal}
        Instances travel from root to leaf; at each node only the \alert{current rule} matters —
        no memory of the path so far.
      \end{alertblock}
  \end{columns}
\end{frame}

\begin{frame}{Univariate Splits and Interpretability}
  \begin{itemize}
    \item Each internal node is typically associated with a partition induced by a
          \alert{single attribute} (univariate splitting).
    \item This constraint is not mathematically necessary, but it plays a crucial role in
          making the model \alert{interpretable}.
    \item Example: \texttt{petal length $\leq$ 2.45} is immediately understandable;
          a combination of features would require further explanation.
    \item As the tree grows deeper, each path from root to leaf accumulates a
          \alert{conjunction of simple conditions}.
  \end{itemize}
  \begin{block}{Leaf as a region}
    A leaf corresponds to the set of all instances satisfying \alert{every condition}
    along the path from the root to that leaf simultaneously — a clearly delineated
    subregion of the feature space.
  \end{block}
\end{frame}

\begin{frame}{Predictions at the Leaves}
  \begin{columns}[t]
    \column{0.48\textwidth}
      \begin{block}{Classification}
        A leaf may store:
        \begin{itemize}
          \item A single \alert{class label} (deterministic prediction).
          \item The \alert{empirical class distribution} observed in that region during training
                (calibrated probability estimates).
        \end{itemize}
      \end{block}
    \column{0.48\textwidth}
      \begin{block}{Regression}
        A leaf stores a \alert{numerical summary} of target values in that region, most
        commonly their \alert{mean}, returned as the predicted output for any new instance
        landing in the leaf.
      \end{block}
  \end{columns}
  \vspace{0.6em}
  \begin{itemize}
    \item The \alert{depth} of the tree and the \alert{number of leaves} reflect the complexity
          of the learned function.
    \item Shallow tree: coarser partition, simpler decision boundary.
    \item Deep tree: finer distinctions, but higher risk of \alert{overfitting}.
  \end{itemize}
\end{frame}

\begin{frame}{Divide-and-Conquer Problem Solving}
  \begin{itemize}
    \item The tree structure encodes \alert{divide-and-conquer} naturally.
    \item The root node faces the hardest problem: predicting the output for the entire,
          presumably heterogeneous, dataset.
    \item By choosing an informative splitting criterion, it divides this hard problem into
          \alert{two (or more) easier sub-problems}, each involving a more homogeneous subset.
    \item Each child node applies the same logic recursively, further reducing heterogeneity.
    \item Recursion terminates when subsets become \alert{sufficiently pure} or a stopping
          criterion is met.
  \end{itemize}
  \begin{alertblock}{Complexity vs.\ generalization}
    A shallow tree imposes a coarser partition; a deep tree can capture finer distinctions
    at the risk of \alert{overfitting} to the training data.
  \end{alertblock}
\end{frame}

%% ════════════════════════════════════════════════════════════════════════════
\section{A Visual Example: The Iris Dataset}

\begin{frame}{The Iris Dataset}
  \framesubtitle{Three species, four features — a classical benchmark}
  \begin{center}
    \includegraphics[width=.65\linewidth]{figures/iris.png}
  \end{center}
  \begin{itemize}
    \item Contains measurements of petal and sepal dimensions for flowers of three species.
    \item Three species form visible clusters in feature space, with some overlap.
    \item A small number of well-chosen splits may already achieve good separation.
  \end{itemize}
\end{frame}

\begin{frame}{After the First Split: Tree}
  \framesubtitle{Initial decision tree with first split}
  \begin{center}
    \includegraphics[width=.55\linewidth]{figures/dec_tree_0.png}
  \end{center}
  \begin{itemize}
    \item A single axis-aligned cut divides instances into \alert{two groups}.
    \item One resulting region may already be pure enough to constitute a leaf.
    \item The other still contains a mixture of classes and must be split further.
  \end{itemize}
\end{frame}

\begin{frame}{After the First Split: Feature Space}
  \framesubtitle{Corresponding partition of the feature space}
  \begin{center}
    \includegraphics[width=.8\linewidth]{figures/dec_tree_regions_0.png}
  \end{center}
  \begin{itemize}
    \item The partition consists of just \alert{two half-planes} separated by a single decision boundary.
    \item Each path root$\to$leaf corresponds to exactly one region in the feature space.
  \end{itemize}
\end{frame}

\begin{frame}{After the Second Split: Tree}
  \framesubtitle{Decision tree after second split}
  \begin{center}
    \includegraphics[width=.55\linewidth]{figures/dec_tree_1.png}
  \end{center}
  \begin{itemize}
    \item The second split refines the partition within the impure region identified at the previous step.
    \item The tree has grown one level deeper in the branch that required further discrimination.
    \item At each step, the algorithm selects the split that \alert{maximally reduces impurity}
          (e.g., Gini index or information gain).
  \end{itemize}
\end{frame}

\begin{frame}{After the Second Split: Feature Space}
  \framesubtitle{Updated feature space partition}
  \begin{center}
    \includegraphics[width=.8\linewidth]{figures/dec_tree_regions_1.png}
  \end{center}
  \begin{itemize}
    \item The feature space is now divided into \alert{three regions}.
    \item Every new split introduces a new axis-aligned boundary within the current node's region.
  \end{itemize}
\end{frame}

\begin{frame}{Final Decision Tree}
  \framesubtitle{Final decision tree after all splits}
  \begin{center}
    \includegraphics[width=.65\linewidth]{figures/dec_tree.png}
  \end{center}
  \begin{itemize}
    \item Each leaf is assigned the class label most common among its training instances.
    \item The boundaries between regions are \alert{rectilinear} — a direct consequence of univariate,
          axis-aligned splits.
  \end{itemize}
\end{frame}

\begin{frame}{Final Partition of the Feature Space}
  \framesubtitle{Final partition of the feature space}
  \begin{center}
    \includegraphics[width=.8\linewidth]{figures/dec_tree_regions.png}
  \end{center}
  \begin{itemize}
    \item \alert{Axis-aligned rectangular regions} distinguish decision trees from linear classifiers
          or SVMs, whose boundaries are oblique hyperplanes.
    \item Every path root$\to$leaf $\Leftrightarrow$ exactly one rectangular cell; the two
          representations are \alert{dual to each other}.
  \end{itemize}
\end{frame}

%% ════════════════════════════════════════════════════════════════════════════
\section{Classification and Construction}

\begin{frame}{Classification: Top-Down Traversal}
  \framesubtitle{Predicting a new instance}
  \begin{itemize}
    \item Let $\mathbf{z} = (z_1, \ldots, z_d)^T$ be a new input instance with $d$ features.
    \item Prediction begins at the \alert{root node}, representing the full feature space.
    \item At each internal node, the model evaluates the node's decision rule and selects the
          outgoing branch corresponding to the outcome.
    \item The instance carries \alert{no memory} of the path already traveled: only the current
          node's rule matters.
  \end{itemize}
  \begin{block}{Threshold test (continuous feature)}
    Suppose the node is associated with feature $x_i$ and threshold $\theta_i$:
    \[
      \text{if } z_i \leq \theta_i \ \to \ \text{go left;} \qquad
      \text{if } z_i > \theta_i \ \to \ \text{go right.}
    \]
    Each such test is geometrically a \alert{cut along an axis-aligned hyperplane}.
  \end{block}
\end{frame}

\begin{frame}{Classification: Output and Efficiency}
  \begin{itemize}
    \item Traversal continues, node by node, until a \alert{leaf} is reached.
    \item At the leaf, the tree returns:
      \begin{itemize}
        \item A \alert{class label} (deterministic prediction), or
        \item A \alert{probability distribution} over classes (from which the most probable class
              can be selected).
      \end{itemize}
  \end{itemize}
  \begin{block}{Computational complexity of inference}
    Because the traversal inspects exactly one node per level of the tree, the time required
    to produce a prediction is \alert{$O(h)$} where $h$ is the height of the tree.
  \end{block}
  \begin{itemize}
    \item Inference is \alert{extremely efficient}: even a large, complex tree requires only a small
          number of simple comparisons.
    \item Contrast with nearest-neighbor classifiers, whose prediction cost \alert{grows with $n$}.
  \end{itemize}
\end{frame}

\begin{frame}{Construction: Greedy Top-Down Algorithm}
  \begin{itemize}
    \item Tree construction proceeds \alert{recursively}, starting at the root.
    \item The most common approach: \alert{greedy and top-down}.
      \begin{itemize}
        \item At each node, choose the locally best split according to some criterion.
        \item Apply it immediately, without reconsidering earlier choices.
        \item Continue independently on each resulting child subset.
      \end{itemize}
    \item The greedy strategy does not search for the \alert{globally optimal tree}
          (computationally intractable in general).
    \item Instead, it builds a reasonable tree efficiently by making \alert{locally good decisions}.
  \end{itemize}
  \begin{alertblock}{Two key questions at every node}
    \textbf{How to partition?} Which feature and threshold yield the most useful separation?\\
    \textbf{When to stop?} When should a region be committed to a prediction rather than split further?
  \end{alertblock}
\end{frame}

\begin{frame}{Construction: Recursive Procedure}
  \begin{enumerate}
    \item Start with the \alert{entire training set} at the root.
    \item For each node associated with subset $S$:
      \begin{itemize}
        \item Decide whether $S$ should become a \alert{leaf} or be split further.
        \item If splitting: select the feature and split rule that best improves purity; partition
              $S$ into child subsets; recurse on each child.
        \item If terminal: assign prediction information — class label, class distribution, or
              numerical value.
      \end{itemize}
  \end{enumerate}
  \vspace{0.4em}
  \begin{block}{Stopping criteria (early stopping)}
    Common conditions for declaring a node a leaf:
    \begin{itemize}
      \item All instances at the node have the same class label (pure node).
      \item The number of instances falls below a minimum threshold.
      \item Maximum depth has been reached.
      \item No split achieves a sufficiently large impurity reduction.
    \end{itemize}
  \end{block}
\end{frame}

%% ════════════════════════════════════════════════════════════════════════════
\section{Impurity Measures and Split Selection}

\begin{frame}{The Notion of Impurity}
  \begin{itemize}
    \item Selection of the best split is guided by the notion of \alert{impurity}.
    \item A node is ``good'' when labels of instances in it are mostly the same
          (\alert{high homogeneity}).
    \item A node is ``bad'' when it contains a mixture of different classes
          (\alert{high heterogeneity}).
    \item Impurity measures formalize this: they assign a real number to a class distribution,
          \alert{small for concentrated (pure) distributions} and \alert{large for uniform ones}.
  \end{itemize}
  \begin{block}{Goal}
    Find the partition that \alert{reduces impurity the most}: move from a more heterogeneous
    parent node to more homogeneous children.
  \end{block}
\end{frame}

\begin{frame}{Formal Definition of Impurity Measures}
  \begin{itemize}
    \item Consider a node with instances from $k$ classes; let $p = (p_1, \ldots, p_k)$ be the
          empirical class distribution ($p_i$ = fraction of class $i$).
    \item An impurity measure is a function $\phi: p \mapsto \mathbb{R}$.
  \end{itemize}
  \begin{block}{Required properties}
    \renewcommand{\arraystretch}{1.3}
    \begin{tabular}{ll}
      \textbf{Non-negativity} & $\phi(p) \geq 0$ for all $p$. \\
      \textbf{Minimum at purity} & $\phi(p) = 0$ if $\exists\, i: p_i = 1$. \\
      \textbf{Maximum at balance} & $\phi(p)$ maximal when $p_i = 1/k\ \forall i$. \\
      \textbf{Symmetry} & $\phi(p) = \phi(p')$ for any permutation $p'$ of $p$. \\
      \textbf{Smoothness} & $\phi$ is differentiable on the interior of the simplex.
    \end{tabular}
  \end{block}
  \begin{itemize}
    \item Intuition: $\phi$ measures how hard it is to \alert{correctly guess the class label}
          of a randomly drawn instance from the node.
  \end{itemize}
\end{frame}

\begin{frame}{Goodness of Split}
  \begin{itemize}
    \item Given a labeled set $S$ with class-probability vector $p_S$, and a split function
          $f: S \to \{1, \ldots, r\}$ producing disjoint subsets:
    \[
      s_i = \{x \in S \mid f(x) = i\}, \qquad S = s_1 \cup \cdots \cup s_r.
    \]
    \item Compare the impurity of the parent with the \alert{expected impurity after the split}:
  \end{itemize}
  \begin{block}{Goodness of split}
    \[
      \Delta_\phi(S, f) = \phi(p_S) - \sum_{i=1}^{r} \frac{|s_i|}{|S|}\,\phi(p_{s_i}).
    \]
    A \alert{larger} $\Delta_\phi$ means the children are, on average, substantially purer than
    the parent. The algorithm selects the split achieving the \alert{highest goodness of split}.
  \end{block}
  \begin{itemize}
    \item The weighting by $|s_i|/|S|$ ensures that \alert{large child nodes} contribute more than
          small ones to the expected post-split impurity.
  \end{itemize}
\end{frame}

\begin{frame}{Practical Split Functions: Discrete Features}
  \begin{itemize}
    \item Practical trees restrict attention to splits based on a \alert{single feature},
          for interpretability and computational feasibility.
  \end{itemize}
  \begin{block}{Discrete (categorical) feature}
    \begin{itemize}
      \item \textbf{Multiway split}: one branch per distinct value of the feature.
      \item \textbf{Binary split}: group the values into two sets.
      \item Binary version preferred for high-cardinality features: multiway splits
            produce many small, statistically unreliable children.
    \end{itemize}
  \end{block}
  \vspace{0.4em}
  \begin{itemize}
    \item For small cardinality, the best binary partition can be found by exhaustive enumeration.
    \item The goal is always the same: maximize $\Delta_\phi$ over all candidate split functions.
  \end{itemize}
\end{frame}

\begin{frame}{Practical Split Functions: Continuous Features}
  \begin{block}{Continuous feature — threshold split}
    For a continuous feature $x_i$ and threshold $\theta$:
    \[
      x_i \leq \theta \quad \text{vs} \quad x_i > \theta.
    \]
    Candidate thresholds are typically chosen from the \alert{set of observed feature values}.
  \end{block}
  \vspace{0.4em}
  \begin{block}{Efficient threshold search}
    \begin{enumerate}
      \item Sort training instances by the feature value: $O(n \log n)$.
      \item Consider all midpoints between consecutive distinct values as candidates.
      \item Evaluate $\Delta_\phi$ at each candidate.
      \item Select the threshold achieving the highest goodness of split.
    \end{enumerate}
    Total cost for a single feature: $O(n \log n)$.
  \end{block}
\end{frame}

\begin{frame}{Entropy as an Impurity Measure}
  \begin{itemize}
    \item A theoretically well-motivated impurity measure from \alert{information theory}: entropy.
    \item Entropy quantifies the \alert{uncertainty} present in a probability distribution.
    \item When applied to the class distribution at a node, it captures how unpredictable the
          class label of an instance reaching that node is.
  \end{itemize}
  \begin{block}{Definition}
    For a set $S$ of instances from $k$ classes, with $S_i$ the subset of class $i$:
    \[
      H_S = -\sum_{i=1}^{k} \frac{|S_i|}{|S|} \log_2 \frac{|S_i|}{|S|}.
    \]
    The base-2 logarithm gives entropy the unit of \alert{bits}: $H_S$ equals the expected number
    of binary questions needed (under an optimal strategy) to identify the class of a
    randomly chosen instance.
  \end{block}
\end{frame}

\begin{frame}{Entropy: Extreme Values and Interpretation}
  \begin{itemize}
    \item \textbf{Pure node} (all instances of class $i$): $|S_i|/|S| = 1$, $|S_j|/|S| = 0$
          for $j \neq i$.
          \[
            H_S = 0. \quad \text{No uncertainty, no information needed.}
          \]
    \item \textbf{Maximally mixed node} ($k$ classes equally represented: $|S_i|/|S| = 1/k$):
          \[
            H_S = \log_2 k. \quad \text{Maximum uncertainty.}
          \]
    \item The upper bound grows with $k$: more classes means more potential uncertainty.
  \end{itemize}
  \begin{alertblock}{Convention}
    By convention, $0 \log_2 0 = 0$, so entropy is well defined even for empty classes.
  \end{alertblock}
\end{frame}

\begin{frame}{Information Gain}
  \begin{itemize}
    \item When entropy is used as the impurity function $\phi$, the goodness-of-split criterion
          $\Delta_\phi(S, f)$ specializes to \alert{information gain}.
    \item For a split function $f$ that partitions $S$ into subsets $s_1, \ldots, s_r$:
  \end{itemize}
  \begin{block}{Information Gain}
    \[
      IG(S, f) = H_S - \sum_{j=1}^{r} \frac{|s_j|}{|S|}\, H_{s_j}.
    \]
    Information gain measures the \alert{reduction in entropy} — the reduction in uncertainty about
    the class label — achieved by applying split $f$ to set $S$.
  \end{block}
  \begin{itemize}
    \item Information gain is the criterion used by the classic \alert{ID3} algorithm.
    \item A known bias: raw information gain favors features with \alert{many distinct values}
          (high cardinality), which may not be truly informative.
    \item \alert{Gain ratio} (used by C4.5) corrects for this bias by normalizing $IG$ by the
          entropy of the split itself.
  \end{itemize}
\end{frame}

\begin{frame}{Gini Impurity}
  \begin{itemize}
    \item Another widely used impurity measure: the \alert{Gini index} (or Gini impurity).
    \item Preferred by the CART framework as the default splitting criterion.
  \end{itemize}
  \begin{block}{Definition}
    For a set $S$ with empirical class distribution $(p_1, \ldots, p_k)$:
    \[
      \text{Gini}(S) = 1 - \sum_{i=1}^{k} p_i^2.
    \]
    Interpretation: the probability that a randomly chosen instance is \alert{misclassified}
    if its label is randomly drawn from the empirical distribution.
  \end{block}
  \begin{columns}[t]
    \column{0.48\textwidth}
      \begin{block}{Gini = 0}
        All instances of one class: $p_i = 1$ for some $i$.
      \end{block}
    \column{0.48\textwidth}
      \begin{block}{Gini maximal}
        All classes equally represented: $\text{Gini} = 1 - 1/k$.
      \end{block}
  \end{columns}
\end{frame}

\begin{frame}{Comparing Impurity Measures}
  \renewcommand{\arraystretch}{1.5}
  \begin{center}
  \begin{tabular}{lll}
    \toprule
    \textbf{Measure} & \textbf{Formula} & \textbf{Used by} \\
    \midrule
    Entropy & $-\sum_i p_i \log_2 p_i$ & ID3, C4.5 \\
    Gain Ratio & $IG(S,f)\,/\,H_f(S)$ & C4.5 \\
    Gini index & $1 - \sum_i p_i^2$ & CART \\
    Variance & $\frac{1}{|S|}\sum(y_i - \bar y)^2$ & Regression trees \\
    \bottomrule
  \end{tabular}
  \end{center}
  \vspace{0.4em}
  \begin{itemize}
    \item Entropy and Gini are \alert{qualitatively similar}: both are symmetric, zero at purity,
          maximal at balance, and smooth on the simplex interior.
    \item In practice, the choice of impurity measure rarely has a large effect on the
          final tree quality; the \alert{stopping and pruning} strategy matters more.
  \end{itemize}
\end{frame}

%% ════════════════════════════════════════════════════════════════════════════
\section{Pruning: Motivation and Overview}

\begin{frame}{The Overfitting Problem}
  \begin{itemize}
    \item Trees can become very expressive by growing deep and creating fine-grained partitions.
    \item This expressiveness comes with a high risk of \alert{overfitting}.
    \item A decision tree allowed to grow without restraint will eventually create \alert{one leaf
          per training instance}: zero training error, but nothing transferable to new data.
    \item Unlike models with continuous parameters (where regularization is smooth), a
          decision tree's complexity emerges from \alert{discrete splitting decisions}.
  \end{itemize}
  \begin{alertblock}{The bias--variance trade-off in trees}
    Managing bias and variance is both \alert{essential} and \alert{genuinely difficult} in
    decision tree models: complexity is controlled by discrete choices rather than by tuning
    a continuous regularization parameter.
  \end{alertblock}
\end{frame}

\begin{frame}{Two Extreme Strategies}
  \begin{columns}[t]
    \column{0.48\textwidth}
      \begin{block}{Early stopping}
        Enforce strict stopping criteria during growth.
        \begin{itemize}
          \item Small, shallow, interpretable trees.
          \item Statistically stable across samples.
          \item \alert{High bias}: may miss deeper patterns.
          \item ``Horizon effect'': a locally poor split might unlock highly informative
                splits further down.
        \end{itemize}
      \end{block}
    \column{0.48\textwidth}
      \begin{alertblock}{Late stopping}
        Loose stopping criteria; tree keeps splitting.
        \begin{itemize}
          \item Training error approaches zero.
          \item Learned distinctions often reflect \alert{noise}.
          \item High variance: sensitive to small fluctuations.
          \item Large, uninterpretable trees.
        \end{itemize}
      \end{alertblock}
  \end{columns}
\end{frame}

\begin{frame}{Pruning: Grow First, Simplify Later}
  \begin{itemize}
    \item Pruning addresses the dilemma with a pragmatic \alert{``grow first, simplify later''} philosophy.
    \item Rather than guessing the correct stopping time during growth, \alert{intentionally grow a
          large tree} that can capture all potentially useful regularities.
    \item Then \alert{systematically remove} branches that do not contribute to generalization.
  \end{itemize}
  \begin{block}{Two phases}
    \textbf{Growing phase}: expand the tree using loose stopping criteria, favoring coverage
    and flexibility.\\[0.3em]
    \textbf{Pruning phase}: delete subtrees that appear to model noise rather than robust
    structure, retaining only parts that improve expected performance on unseen data.
  \end{block}
  \begin{itemize}
    \item Generalization can be estimated via: a held-out \alert{validation set}, a
          \alert{statistical correction} to training error, or a \alert{regularization term}
          penalizing model complexity.
  \end{itemize}
\end{frame}

%% ════════════════════════════════════════════════════════════════════════════
\section{Pruning Methods}

\begin{frame}{Reduced Error Pruning (REP)}
  \framesubtitle{Conceptually the simplest pruning method}
  \begin{itemize}
    \item Data partitioned into three sets: \alert{training set}, \alert{pruning set} (validation),
          and optional test set.
    \item All pruning decisions are made by directly measuring their effect on accuracy on
          the held-out \alert{pruning set}.
  \end{itemize}
  \begin{block}{Algorithm (bottom-up)}
    At each internal node, ask: what if this node were replaced by a leaf predicting the
    majority class among training instances that reached it?\\
    $\to$ Evaluate both the original subtree and the candidate leaf on the pruning set.\\
    $\to$ Replace the subtree with the leaf if doing so \alert{does not decrease accuracy} on
    the pruning set.\\
    $\to$ Repeat until no further pruning improves or maintains accuracy.
  \end{block}
\end{frame}

\begin{frame}{REP: Properties and Limitations}
  \begin{columns}[t]
    \column{0.48\textwidth}
      \begin{block}{Strengths}
        \begin{itemize}
          \item \alert{Exact}: directly optimizes validation accuracy.
          \item \alert{Efficient}: a single bottom-up pass suffices.
          \item Produces small trees that generalize well when the pruning set is large.
        \end{itemize}
      \end{block}
    \column{0.48\textwidth}
      \begin{alertblock}{Limitation}
        \begin{itemize}
          \item Depends on the \alert{size of the pruning set}.
          \item Scarce data: allocating data to pruning leaves too little for growing.
          \item Small pruning set: noisy validation accuracy estimates.
          \item Most attractive when \alert{large amounts of data} are available.
        \end{itemize}
      \end{alertblock}
  \end{columns}
\end{frame}

\begin{frame}{Cost Complexity Pruning (CCP) — CART Framework}
  \framesubtitle{Penalized error criterion balancing accuracy against tree size}
  \begin{itemize}
    \item CCP (also called \alert{weakest-link pruning}) is the strategy used in CART.
    \item Instead of a separate validation set, employs a \alert{penalized error criterion}.
    \item Generates an entire sequence of progressively simpler trees.
  \end{itemize}
  \begin{block}{Cost-complexity measure}
    For any subtree $T$ of the fully grown tree $T_{\max}$:
    \[
      R_\alpha(T) = R(T) + \alpha\,|\tilde{T}|,
    \]
    where $R(T)$ is the resubstitution error, $|\tilde{T}|$ is the number of leaf nodes, and
    $\alpha \geq 0$ is the \alert{complexity parameter} (penalty per leaf).
  \end{block}
  \begin{itemize}
    \item $\alpha = 0$: no penalty $\Rightarrow$ $T_{\max}$ is optimal.
    \item $\alpha$ large: smaller trees preferred (each leaf must justify itself by reducing error by at
          least $\alpha$).
  \end{itemize}
\end{frame}

\begin{frame}{CCP: The Nested Sequence of Optimal Trees}
  \begin{itemize}
    \item For a given $\alpha$, the optimal subtree $T(\alpha)$ is the smallest subtree of
          $T_{\max}$ that minimizes $R_\alpha$.
    \item A fundamental result: optimal subtrees form a finite \alert{nested sequence}:
    \[
      T_{\max} = T_0 \supset T_1 \supset T_2 \supset \cdots \supset T_m = \{t_1\}.
    \]
    \item Computed by repeatedly identifying and collapsing the \alert{weakest link}: the internal
          node $t$ whose removal produces the smallest increase in resubstitution error per leaf removed.
    \item For each internal node $t$, the \alert{effective complexity parameter}:
    \[
      \alpha(t) = \frac{R(t) - R(T_t)}{|\tilde{T}_t| - 1},
    \]
    where $T_t$ is the subtree rooted at $t$ and $|\tilde{T}_t|$ its number of leaves.
  \end{itemize}
\end{frame}

\begin{frame}{CCP: Model Selection via Cross-Validation}
  \begin{itemize}
    \item CCP \alert{avoids specifying $\alpha$ in advance}: the algorithm generates all candidate
          trees in the nested sequence.
    \item \alert{Cross-validation} is then used to select the best value of $\alpha$ (and hence the
          best tree in the sequence).
    \item This makes CCP both \alert{theoretically principled} and \alert{practically flexible}.
  \end{itemize}
  \begin{block}{CART design choices}
    CART builds \alert{binary trees}: every internal node has exactly two children, produced by a
    binary split on a single feature. This:
    \begin{itemize}
      \item Simplifies tree structure and pruning algorithm.
      \item Facilitates unified handling of categorical and continuous features.
      \item For categorical features with $m$ values: considers all $2^{m-1}-1$ binary partitions.
    \end{itemize}
  \end{block}
  \begin{itemize}
    \item CCP with cross-validation is the \alert{default strategy in scikit-learn}.
  \end{itemize}
\end{frame}

\begin{frame}{Pessimistic Error Pruning (PEP) — C4.5 Framework}
  \framesubtitle{Estimating generalization from training data alone}
  \begin{itemize}
    \item PEP is the pruning strategy used in the \alert{C4.5} algorithm (extension of ID3).
    \item Unlike REP, PEP \alert{does not require a separate validation set}: estimates generalization
          error directly from training data using a \alert{statistical upper confidence bound}.
    \item Particularly attractive in settings where available data are \alert{limited}.
  \end{itemize}
  \begin{block}{Motivation}
    The resubstitution error of a leaf is \alert{optimistically biased}: a leaf with zero training
    errors does not necessarily generalize perfectly — it may have seen too few instances for its
    error rate to be estimated reliably.
  \end{block}
\end{frame}

\begin{frame}{PEP: Pessimistic Estimate}
  \begin{itemize}
    \item Suppose a leaf contains $n$ instances and misclassifies $e$ of them.
    \item Rather than using $e/n$, C4.5 treats the observed error rate as a binomial proportion
          and computes an \alert{upper confidence bound} $\bar{e}(n,e)$ at a user-specified
          confidence level (typically $25\%$ — deliberately conservative to encourage pruning).
    \item Estimated error of the leaf: $n \cdot \bar{e}(n, e)$.
    \item For any subtree rooted at an internal node: sum of pessimistic errors of all its leaves.
  \end{itemize}
  \begin{block}{Pruning rule (top-down)}
    Replace a subtree by a leaf if the \alert{estimated error of the leaf} does not exceed the
    sum of estimated errors of the subtree's leaves.\\
    Upper confidence bounds penalize \alert{small leaf sizes}: a leaf with few instances has a
    wide confidence interval, discouraging very fine-grained splits supported by little data.
  \end{block}
\end{frame}

\begin{frame}{C4.5: Additional Design Features}
  \begin{columns}[t]
    \column{0.48\textwidth}
      \begin{block}{Splitting}
        \begin{itemize}
          \item \alert{Multiway splits} for categorical features (one branch per value).
          \item \alert{Gain ratio} as splitting criterion, correcting the bias of raw information
                gain toward high-cardinality features.
        \end{itemize}
      \end{block}
    \column{0.48\textwidth}
      \begin{block}{Missing values}
        \alert{Fractional instance mechanism}: when a training instance has a missing value for
        the splitting feature, it is distributed fractionally among all branches in proportion
        to the fraction of known instances following each branch.
      \end{block}
  \end{columns}
  \vspace{0.4em}
  \begin{itemize}
    \item \alert{C5.0} (commercial successor): introduced boosting integration and more efficient
          memory use, retaining the same fundamental pruning philosophy.
  \end{itemize}
\end{frame}

\begin{frame}{Minimum Error Pruning (MEP)}
  \framesubtitle{A Bayesian perspective on pruning}
  \begin{itemize}
    \item MEP estimates error through a \alert{prior distribution} over class probabilities rather
          than through confidence bounds or a validation set.
  \end{itemize}
  \begin{block}{m-estimate (Laplace-corrected estimate)}
    For a node with $n$ instances, $n_c$ in the majority class, smoothing parameter $m$,
    and prior $p_c = 1/k$:
    \[
      \hat{p}_c = \frac{n_c + m \cdot p_c}{n + m}.
    \]
    Expected error rate: $1 - \hat{p}_c$.
    \begin{itemize}
      \item Small $n$: estimate pulled toward the prior $1/k$ (high uncertainty).
      \item Large $n$: approaches the empirical error rate $1 - n_c/n$.
    \end{itemize}
  \end{block}
\end{frame}

\begin{frame}{MEP: Procedure, Strengths, and Limitations}
  \begin{block}{Pruning rule (bottom-up)}
    Replace a subtree by a leaf if the \alert{expected error of the leaf} is no greater than
    the average expected error of the subtree (weighted by instances reaching each leaf).\\
    $\Rightarrow$ Prunes subtrees supported by too few instances to reliably improve upon
    the prior estimate.
  \end{block}
  \begin{columns}[t]
    \column{0.48\textwidth}
      \begin{block}{Strengths}
        \begin{itemize}
          \item No dedicated validation set required.
          \item No complex statistical bounds.
          \item Efficient: single bottom-up pass.
        \end{itemize}
      \end{block}
    \column{0.48\textwidth}
      \begin{alertblock}{Limitation}
        Sensitivity to the choice of smoothing constant $m$: different values lead to
        substantially different pruning levels; \alert{no universally agreed method} for
        choosing $m$ automatically.
      \end{alertblock}
  \end{columns}
\end{frame}

\begin{frame}{Error-Based Pruning (EBP)}
  \begin{itemize}
    \item EBP is conceptually related to PEP but applies the confidence-bound correction
          \alert{subtree by subtree} rather than leaf by leaf.
    \item The estimated error of a subtree is computed by \alert{aggregating} the confidence-bound-corrected
          errors of its individual leaves.
    \item This is compared against the confidence-bound-corrected error of the corresponding leaf
          that would result from pruning.
    \item The subtree is pruned if the leaf's estimated error is \alert{not substantially worse}
          than the subtree's.
  \end{itemize}
  \begin{block}{Comparison with PEP}
    EBP accounts more carefully for error at the \alert{subtree level}, rather than approximating the
    subtree error as a simple sum of leaf errors. In practice, EBP tends to produce trees of
    \alert{intermediate size} between REP and PEP.
  \end{block}
\end{frame}

\begin{frame}{Minimum Description Length (MDL) Pruning}
  \framesubtitle{An information-theoretic perspective}
  \begin{itemize}
    \item The \alert{MDL principle}: the best model minimizes the total number of bits needed to
          describe both the model and the data given the model.
    \item For decision trees: trade-off between the \alert{description cost of the tree} (grows
          with nodes) and the \alert{description cost of the training labels given the tree}
          (decreases as the tree fits data better).
  \end{itemize}
  \begin{block}{MDL score}
    \[
      MDL(T, S) = L(T) + L(S \mid T),
    \]
    where $L(T)$ is the code length to describe the tree structure and split rules, and
    $L(S \mid T)$ is the code length to describe the training labels given the tree.\\
    $\Rightarrow$ Prune if doing so \alert{reduces the total MDL score}.
  \end{block}
\end{frame}

\begin{frame}{MDL Pruning: Theoretical Appeal and Practical Status}
  \begin{columns}[t]
    \column{0.48\textwidth}
      \begin{block}{Theoretical appeal}
        \begin{itemize}
          \item Coherent framework grounded in \alert{information theory}.
          \item No arbitrary confidence parameter.
          \item Pruning as \alert{model selection}: the simplest model that still explains the data.
        \end{itemize}
      \end{block}
    \column{0.48\textwidth}
      \begin{alertblock}{Practical limitations}
        \begin{itemize}
          \item Requires careful specification of \alert{coding schemes}.
          \item Different coding choices can lead to different pruning decisions.
          \item Less commonly used in off-the-shelf implementations than REP or CCP.
        \end{itemize}
      \end{alertblock}
  \end{columns}
  \vspace{0.4em}
  \begin{itemize}
    \item MDL pruning remains an important conceptual tool for understanding pruning as
          \alert{principled model selection}.
  \end{itemize}
\end{frame}

%% ════════════════════════════════════════════════════════════════════════════
\section{Comparison and Practical Guidelines}

\begin{frame}{Pruning Methods: A Practical Comparison}
  \renewcommand{\arraystretch}{1.4}
  \begin{center}
  \begin{tabular}{lll}
    \toprule
    \textbf{Method} & \textbf{Data required} & \textbf{Key parameter} \\
    \midrule
    REP  & Validation set & (none) \\
    CCP (CART) & Cross-validation & $\alpha$ (auto via CV) \\
    PEP (C4.5) & Training only & Confidence level \\
    EBP  & Training only & Confidence level \\
    MEP  & Training only & Smoothing $m$ \\
    MDL  & Training only & Coding scheme \\
    \bottomrule
  \end{tabular}
  \end{center}
\end{frame}

\begin{frame}{Major Implementations: Structural Differences}
  \begin{columns}[t]
    \column{0.48\textwidth}
      \begin{block}{CART}
        \begin{itemize}
          \item \alert{Binary trees} only.
          \item Unified regression and classification.
          \item \alert{Gini impurity} as default.
          \item Cost-sensitive via weighted misclassification costs.
        \end{itemize}
      \end{block}
      \vspace{0.4em}
      \begin{block}{ID3}
        \begin{itemize}
          \item Raw \alert{information gain}.
          \item Does not handle continuous features or missing values natively.
          \item Historical ancestor of C4.5.
        \end{itemize}
      \end{block}
    \column{0.48\textwidth}
      \begin{block}{C4.5}
        \begin{itemize}
          \item \alert{Multiway splits} for categorical features.
          \item \alert{Gain ratio} to correct high-cardinality bias.
          \item Fractional instance propagation for \alert{missing values}.
          \item PEP pruning; no validation set needed.
        \end{itemize}
      \end{block}
      \vspace{0.4em}
      \begin{block}{Modern (LightGBM, XGBoost)}
        \begin{itemize}
          \item Leaf-wise growth strategies.
          \item Histogram-based split finding.
          \item Explicit regularization in splitting criterion.
        \end{itemize}
      \end{block}
  \end{columns}
\end{frame}

\begin{frame}{Practical Guidelines for Pruning Strategy}
  \begin{itemize}
    \item An \alert{unpruned tree almost always overfits}; a reasonably pruned tree typically
          achieves substantially better generalization regardless of the specific method used.
  \end{itemize}
  \begin{block}{Guidelines by setting}
    \begin{tabular}{ll}
      \toprule
      \textbf{Setting} & \textbf{Recommended approach} \\
      \midrule
      Large data & REP or CCP with cross-validation \\
      Scarce data & PEP or MEP (no validation set needed) \\
      Theoretical interpretability & MDL pruning \\
      Decision tree in an ensemble & Limit maximum depth directly \\
      \bottomrule
    \end{tabular}
  \end{block}
  \begin{itemize}
    \item When the decision tree is used within an \alert{ensemble} (random forests, gradient
          boosting), the ensemble provides its own variance reduction — aggressive pruning at the
          individual tree level is often \alert{unnecessary}.
  \end{itemize}
\end{frame}

%% ════════════════════════════════════════════════════════════════════════════
\section{Summary}

\begin{frame}{Summary: Key Concepts}
  \renewcommand{\arraystretch}{1.5}
  \begin{center}
  \begin{tabular}{ll}
    \toprule
    \textbf{Concept} & \textbf{Definition / Formula} \\
    \midrule
    Impurity (general) & $\phi(p) \geq 0$, zero at purity, max at balance \\
    Goodness of split & $\Delta_\phi(S,f) = \phi(p_S) - \sum_i \frac{|s_i|}{|S|}\phi(p_{s_i})$ \\
    Entropy & $H_S = -\sum_i \frac{|S_i|}{|S|}\log_2\frac{|S_i|}{|S|}$ \\
    Information gain & $IG(S,f) = H_S - \sum_j \frac{|s_j|}{|S|}H_{s_j}$ \\
    Gini index & $\text{Gini}(S) = 1 - \sum_i p_i^2$ \\
    CCP penalty & $R_\alpha(T) = R(T) + \alpha|\tilde{T}|$ \\
    Weakest link & $\alpha(t) = (R(t)-R(T_t))/(|\tilde{T}_t|-1)$ \\
    m-estimate (MEP) & $\hat{p}_c = (n_c + m/k)/(n+m)$ \\
    MDL score & $MDL(T,S) = L(T) + L(S|T)$ \\
    Inference cost & $O(h)$, $h$ = height of the tree \\
    \bottomrule
  \end{tabular}
  \end{center}
\end{frame}

\begin{frame}{Summary: Methods at a Glance}
  \begin{columns}[t]
    \column{0.48\textwidth}
      \begin{block}{Tree Structure}
        \begin{itemize}
          \item Root, internal nodes, leaves.
          \item Univariate splits $\Rightarrow$ axis-aligned boundaries.
          \item Leaves: class label, distribution, or mean.
          \item Depth $\leftrightarrow$ complexity.
        \end{itemize}
      \end{block}
      \vspace{0.4em}
      \begin{block}{Impurity \& Splitting}
        \begin{itemize}
          \item Entropy (ID3, C4.5) / Gain ratio (C4.5).
          \item Gini (CART).
          \item Threshold search: $O(n\log n)$ per feature.
          \item Best split = max $\Delta_\phi$.
        \end{itemize}
      \end{block}
    \column{0.48\textwidth}
      \begin{block}{Pruning}
        \begin{itemize}
          \item REP: exact, needs validation set.
          \item CCP: nested sequence + CV (scikit-learn default).
          \item PEP/EBP: training-only, C4.5.
          \item MEP: Bayesian, smoothing $m$.
          \item MDL: information-theoretic.
        \end{itemize}
      \end{block}
      \vspace{0.4em}
      \begin{block}{Practical notes}
        \begin{itemize}
          \item Always apply some pruning.
          \item In ensembles: limit depth instead.
          \item Interpretability $\propto$ shallowness.
        \end{itemize}
      \end{block}
  \end{columns}
\end{frame}

%\end{document}
